\documentclass[12pt,a4paper]{article}

% ── Packages ──────────────────────────────────────────────────────
\usepackage[utf8]{inputenc}
\usepackage[T1]{fontenc}
\usepackage{amsmath,amssymb,amsthm}
\usepackage{mathtools}
\usepackage{physics}
\usepackage{graphicx}
\usepackage[margin=1in]{geometry}
\usepackage{hyperref}
\usepackage{cleveref}
\usepackage{booktabs}
\usepackage{multirow}
\usepackage{xcolor}
\usepackage{float}
\usepackage{enumitem}

% ── Theorem environments ──────────────────────────────────────────
\newtheorem{definition}{Definition}
\newtheorem{conjecture}{Conjecture}
\newtheorem{proposition}{Proposition}
\newtheorem{remark}{Remark}

% ── Shortcuts ─────────────────────────────────────────────────────
\newcommand{\SvN}{S_{\mathrm{vN}}}
\newcommand{\IvN}{I_{\mathrm{vN}}}
\newcommand{\rhoQ}{\rho_Q}
\newcommand{\muQ}{\mu_Q}
\newcommand{\calL}{\mathcal{L}}
\newcommand{\calH}{\mathcal{H}}
\renewcommand{\Tr}{\mathrm{Tr}}
\newcommand{\nats}{\,\mathrm{nats}}

% ══════════════════════════════════════════════════════════════════
\begin{document}

\title{Falsification of the Born--Fisher--Experiential Conjecture\\
in a Qubit Toy Model}

\author{Bryan Ehrlich}

\date{March 2026}

\maketitle

% ── Abstract ──────────────────────────────────────────────────────
\begin{abstract}
The experiential density $\rhoQ = \IvN(B;M)(1 - \IvN/\SvN(B))$
measures meaningful complexity in bipartite quantum systems by
penalising both under-correlated and over-correlated regimes.
The Born--Fisher--Experiential conjecture proposes that Born-rule
probability assignments are dynamically selected by the
half-saturation condition $\IvN = \SvN(B)/2$, where the
experiential density is maximal.
We test this conjecture numerically using a qubit toy model with two
complementary approaches: (i) a static diagonal-state sweep over a
$200 \times 200$ grid in the body-probability / tracking-accuracy
parameter space, and (ii) Lindblad master equation dynamics with an
exchange Hamiltonian and dephasing.
The static test shows that the half-saturation tracking accuracy
$\alpha_{1/2}(p)$ varies with $p$ (spread $0.069$) but Born-rule
distributions display no special property compared with non-Born
alternatives.
The dynamical test reveals that $\rhoQ(t) \le 0$ throughout all
$1\,900^{+}$ Lindblad trajectories because $\IvN/\SvN(B)$
monotonically decreases from $2$ (pure entangled) to $1$ (perfect
tracking), never entering the regime where $\rhoQ$ is positive.
The trajectory functional $\muQ(\theta)$ is identically zero for all
initial states and all tested decoherence parameters.
The conjecture is cleanly falsified for this model class.
We identify the exchange Hamiltonian's perfect-tracking behaviour as
the mechanism and suggest non-Markovian channels as a potential avenue
for modified conjectures.
\end{abstract}

% ══════════════════════════════════════════════════════════════════
\section{Introduction}\label{sec:intro}

A central question in the foundations of quantum mechanics is why the
Born rule holds: given a quantum state $\ket{\psi} = \sum_k c_k
\ket{k}$, the probability of measuring outcome $k$ is $p_k =
|c_k|^2$.  Despite its empirical success, the rule remains
axiomatic in standard quantum theory.  Several derivation programmes
exist---Gleason's theorem~\cite{Gleason1957}, operational
approaches~\cite{Chiribella2011}, and Valentini's
sub-quantum $H$-theorem~\cite{Valentini2005}.  Decision-theoretic
arguments due to Deutsch~\cite{Deutsch1999} and
Wallace~\cite{Wallace2007} derive Born-rule probabilities from
rationality axioms applied to Everettian branching, while Zurek's
envariance programme~\cite{Zurek2005} obtains $p_k = |c_k|^2$ from
entanglement symmetry.  However, none of these approaches employs a
bipartite mutual-information functional as a variational selector
for probability assignments.

The experiential measure framework~\cite{Ehrlich2026a} introduces
a scalar functional on bipartite systems that quantifies meaningful
complexity.  For a classical system with body $B$ and model $M$, the
experiential density is
\begin{equation}\label{eq:rho-classical}
  \rho = I(B;M)\!\left(1 - \frac{I(B;M)}{H(B)}\right),
\end{equation}
where $I(B;M)$ is the mutual information and $H(B)$ is the marginal
entropy.  This functional vanishes at both extremes of correlation
($I = 0$: no model, $I = H$: perfect model) and peaks at half
saturation $I = H/2$, where $\rho_{\max} = H/4$.  The peak
describes a system whose model captures exactly half of the body's
entropy---a regime of maximal ``surprise-generating capacity''.

The quantum extension~\cite{Ehrlich2026a} replaces Shannon
entropies with von Neumann entropies:
\begin{definition}[Quantum experiential density]\label{def:rhoQ}
  For a bipartite density matrix $\rho_{BM}$ with reduced states
  $\rho_B = \Tr_M(\rho_{BM})$ and $\rho_M = \Tr_B(\rho_{BM})$,
  define
  \begin{equation}\label{eq:rhoQ}
    \rhoQ \;=\; \IvN(B;M)\!\left(
      1 - \frac{\IvN(B;M)}{\SvN(B)}
    \right),
  \end{equation}
  where $\SvN(\rho) = -\Tr(\rho \ln \rho)$ is the von Neumann
  entropy and $\IvN(B;M) = \SvN(B) + \SvN(M) - \SvN(BM)$ is the
  quantum mutual information.
\end{definition}

A critical constraint shapes the entire analysis.  For any pure
bipartite state $\ket{\psi}_{BM}$, the Schmidt decomposition gives
$\SvN(BM) = 0$ and $\SvN(B) = \SvN(M)$, so
\begin{equation}\label{eq:pure-ratio}
  \IvN(B;M) = 2\,\SvN(B),
  \qquad
  \rhoQ = 2\,\SvN(B)(1 - 2) = -2\,\SvN(B) < 0.
\end{equation}
The quantum experiential density is \emph{negative} for all entangled
pure states because quantum mutual information can exceed the
marginal entropy---a phenomenon with no classical analogue.  The
functional $\rhoQ$ is therefore physically meaningful only in the
decohered (mixed-state) regime where $\IvN \le \SvN(B)$.

\begin{conjecture}[Born--Fisher--Experiential]\label{conj:BFE}
  Let $\ket{\psi} = \sum_k c_k \ket{k}$ be a quantum state
  measured in the pointer basis $\{\ket{k}\}$.
  Among all probability assignments $\{p_k\}$ over measurement
  outcomes, the Born-rule assignment $p_k = |c_k|^2$ is
  dynamically selected by the condition that the time-integrated
  experiential density
  \begin{equation}\label{eq:muQ}
    \muQ(\theta)
    \;=\; \int_0^{T} \max\!\bigl(\rhoQ(t),\, 0\bigr)\, dt
  \end{equation}
  is extremised at the Born-rule initial state, where $\theta$
  parametrises the initial entanglement angle.
  Equivalently, the half-saturation condition
  $\IvN = \SvN(B)/2$ plays a distinguished role for Born-rule
  probability assignments.
\end{conjecture}

The conjecture connects the observer-dependent experiential measure
to quantum foundations: if true, it would provide a variational
mechanism for why nature selects $p_k = |c_k|^2$ rather than some
other probability rule.  The connection to Fisher information arises
through the Baez--Dolan groupoid cardinality~\cite{Baez2001}, which
motivates the $1/|\mathrm{Aut}(x)|$ weighting that appears
naturally in structure spaces.  (The Fisher information connection is
motivational; the detailed construction appears in the companion
work~\cite{Ehrlich2026a} and is not computed in the present paper.)

In this paper we test \cref{conj:BFE} using a qubit toy model
($d = 2$) with two complementary numerical strategies.  Test~A
(\cref{sec:testA}) performs a static sweep over diagonal
body--model states, asking whether the half-saturation tracking
accuracy depends on whether the body probability follows the Born
rule.  Test~B (\cref{sec:testB}) integrates the full Lindblad master
equation starting from pure entangled states, tracking $\rhoQ(t)$
and the trajectory functional $\muQ(\theta)$ across $1\,900^{+}$
parameter combinations.

Our results are unambiguous.  Test~A shows that $\alpha_{1/2}(p)$
varies with $p$ but Born-rule distributions possess no special
property at half saturation.  Test~B shows that $\rhoQ(t) \le 0$
throughout the entire Lindblad evolution for every tested trajectory,
making $\muQ \equiv 0$ with no dynamical selection mechanism.  The
conjecture is falsified for this class of Lindblad models.

% ══════════════════════════════════════════════════════════════════
\section{Model and Methods}\label{sec:methods}

\subsection{Quantum experiential density}

We work in the bipartite Hilbert space $\calH = \calH_B \otimes
\calH_M$ with $\dim \calH_B = \dim \calH_M = d$.  The body
subsystem $B$ represents the physical system being observed; the
model subsystem $M$ represents the observer's internal
representation.  All entropies use the natural logarithm (nats).

The quantities entering the experiential density
\eqref{eq:rhoQ} are computed from the joint density matrix
$\rho_{BM}$ via eigenvalue decomposition:
\begin{equation}\label{eq:SvN}
  \SvN(\rho) = -\sum_i \lambda_i \ln \lambda_i,
\end{equation}
where $\{\lambda_i\}$ are the eigenvalues of $\rho$, with the
convention $0 \ln 0 = 0$.  Eigenvalues below a cutoff
$\varepsilon = 10^{-15}$ are treated as zero.  Reduced density
matrices are obtained by partial trace:
$\rho_B = \Tr_M(\rho_{BM})$, $\rho_M = \Tr_B(\rho_{BM})$.

We define the \emph{information ratio}
\begin{equation}\label{eq:ratio}
  r \;\equiv\; \frac{\IvN(B;M)}{\SvN(B)},
\end{equation}
with the convention $r = 0$ when $\SvN(B) = 0$.  The experiential
density is positive if and only if $0 < r < 1$, i.e.\ the system
is in the under-correlated regime.  For $r > 1$ (entangled/
over-correlated), $\rhoQ < 0$.  For pure states, $r = 2$ always.

\subsection{Diagonal BM state parametrisation}\label{sec:diagonal}

For the static test, we construct diagonal body--model states
directly in the pointer basis.  A qubit body with probability $p$
of being in state $\ket{0}$ and a model with tracking accuracy
$\alpha$ yields the joint state:
\begin{equation}\label{eq:diagonal-state}
  \rho_{BM}(p,\alpha) = p\,\ketbra{0}{0} \otimes \rho_M^{(0)}
  + (1-p)\,\ketbra{1}{1} \otimes \rho_M^{(1)},
\end{equation}
where
\begin{equation}\label{eq:model-states}
  \rho_M^{(0)} = \begin{pmatrix} \alpha & 0 \\ 0 & 1-\alpha
  \end{pmatrix},
  \qquad
  \rho_M^{(1)} = \begin{pmatrix} 1-\alpha & 0 \\ 0 & \alpha
  \end{pmatrix}.
\end{equation}
When $\alpha = 1$ the model perfectly tracks the body
($\IvN = \SvN(B)$, $r = 1$, $\rhoQ = 0$).  When
$\alpha = 1/d = 1/2$ the model carries no information
($\IvN = 0$, $\rhoQ = 0$).  For intermediate $\alpha$, the
experiential density is positive.

For a given $p$, the half-saturation tracking accuracy
$\alpha_{1/2}(p)$ is defined as the value of $\alpha$ at which
$r = 1/2$.  We extract $\alpha_{1/2}$ by bisection root-finding on
$r(\alpha) - 1/2 = 0$, converging to tolerance $10^{-12}$.

The Born rule enters through the parametrisation of the body
probability.  For a quantum state $\ket{\psi} = \cos\theta\ket{0}
+ \sin\theta\ket{1}$, the Born rule gives $p = \cos^2\theta$.
Non-Born alternatives include $p = \cos^4\theta / Z$ (squared Born
rule), $p = |\cos\theta| / Z$ (absolute-value rule), and
$p = 1/2$ (uniform), where $Z$ is a normalisation constant.

\subsection{Lindblad master equation}\label{sec:lindblad}

For the dynamical test, we evolve a pure entangled initial state
\begin{equation}\label{eq:initial}
  \ket{\psi(\theta)} = \cos\theta\,\ket{00} + \sin\theta\,\ket{11}
\end{equation}
under the Lindblad master equation
\begin{equation}\label{eq:lindblad}
  \frac{d\rho}{dt} = -i[H_{\mathrm{int}}, \rho]
  + \sum_k \left(
    L_k \rho L_k^\dagger
    - \tfrac{1}{2}\{L_k^\dagger L_k, \rho\}
  \right).
\end{equation}

The interaction Hamiltonian is the exchange (or swap) coupling,
\begin{equation}\label{eq:exchange}
  H_{\mathrm{int}} = g\,(\sigma_x \otimes \sigma_x
  + \sigma_y \otimes \sigma_y),
\end{equation}
which drives the model subsystem to track the body.  The collapse
operator implements dephasing in the pointer basis:
\begin{equation}\label{eq:collapse}
  L = \sqrt{\gamma_D}\,(\sigma_z \otimes I_2),
\end{equation}
where $\gamma_D$ is the dephasing rate.  This gives a decoherence
time $\tau_D = 1/(2\gamma_D)$.

For the asymmetric extension, we replace the single collapse
operator with two projective operators:
\begin{equation}\label{eq:asym-collapse}
  L_0 = \sqrt{\gamma_0}\,(\ketbra{0}{0} \otimes I_2),
  \qquad
  L_1 = \sqrt{\gamma_1}\,(\ketbra{1}{1} \otimes I_2),
\end{equation}
with $\gamma_0 / \gamma_1 \ne 1$ breaking the symmetry between
pointer states.

\subsection{Superoperator construction and integration}

We vectorise the density matrix by column-stacking,
$\mathrm{vec}(\rho) = \rho.\mathrm{flatten}(\text{`F'})$, so that
$\mathrm{vec}(A\rho B) = (B^T \otimes A)\,\mathrm{vec}(\rho)$.
The Lindblad equation becomes
\begin{equation}\label{eq:super}
  \frac{d}{dt}\mathrm{vec}(\rho) = \calL \cdot \mathrm{vec}(\rho),
\end{equation}
where the superoperator $\calL$ is a $d^4 \times d^4$ matrix
($16 \times 16$ for qubits):
\begin{equation}\label{eq:super-explicit}
  \calL = -i\bigl(I \otimes H - H^T \otimes I\bigr)
  + \sum_k \left(
    L_k^* \otimes L_k
    - \tfrac{1}{2}\, I \otimes L_k^\dagger L_k
    - \tfrac{1}{2}\,(L_k^\dagger L_k)^T \otimes I
  \right).
\end{equation}

We integrate \eqref{eq:super} using a fixed-step fourth-order
Runge--Kutta (RK4) scheme with $n = 2\,000$ steps over the interval
$[0, 10/\gamma_D]$.  Convergence is verified by comparing with
$n = 4\,000$ steps; the difference in $\muQ$ is below $10^{-15}
\nats\!\cdot\!\mathrm{time}$.

At each stored time step, we reconstruct the density matrix,
enforce Hermiticity by averaging $\rho \to (\rho +
\rho^\dagger)/2$, compute $\SvN(B)$, $\IvN(B;M)$, the ratio $r$,
and $\rhoQ$.  The trajectory functional is
\begin{equation}
  \muQ(\theta) = \int_0^{T_{\mathrm{final}}}
  \max\!\bigl(\rhoQ(t),\, 0\bigr)\, dt,
\end{equation}
evaluated by the trapezoidal rule on the stored time grid.

% ══════════════════════════════════════════════════════════════════
\section{Test A: Static Diagonal States}\label{sec:testA}

\subsection{Grid sweep}

We evaluate the information ratio $r = \IvN/\SvN(B)$ and the
experiential density $\rhoQ$ on a $200 \times 200$ grid over
$(p, \alpha) \in [0.01, 0.99] \times [0.51, 0.999]$ for $d = 2$.
The ratio $r$ spans $[1.4 \times 10^{-4},\, 0.989]$ across the
grid, confirming that the diagonal-state parametrisation covers
the full under-correlated regime.  The experiential density spans
$[7.9 \times 10^{-6},\, 0.173]\nats$, with maximum value $\SvN(B)/4
= \ln(2)/4 = 0.173\nats$ at $p = 0.5$, $\alpha \approx 0.890$,
confirming the parabolic peak structure to $0.002\%$.

\subsection{Half-saturation curve}

For each value of $p$ in the grid, we extract $\alpha_{1/2}(p)$ by
bisection on $r(\alpha) - 1/2 = 0$.  The resulting curve is
\emph{not} constant: $\alpha_{1/2}$ ranges from $0.890$ (at
$p = 0.5$) to $0.959$ (near $p \to 0$ or $p \to 1$), with a
spread of $0.069$ and standard deviation $0.016$.  The curve is
symmetric about $p = 0.5$, as expected from the symmetry
$I(p, \alpha) = I(1-p, \alpha)$.

Representative values:
\begin{center}
\begin{tabular}{lcccccc}
  \toprule
  $p$ & 0.10 & 0.20 & 0.30 & 0.50 & 0.80 & 0.90 \\
  \midrule
  $\alpha_{1/2}$ & 0.916 & 0.901 & 0.894 & 0.890 & 0.901 & 0.916 \\
  \bottomrule
\end{tabular}
\end{center}

\subsection{Born vs.\ non-Born comparison}

We parametrise the body probability $p$ as a function of the
quantum state angle $\theta \in (0, \pi/2)$ using four rules:
\begin{enumerate}[label=(\roman*)]
  \item Born rule: $p = \cos^2\theta$;
  \item Squared Born: $p = \cos^4\theta / Z$;
  \item Absolute-value: $p = |\cos\theta| / Z$;
  \item Uniform: $p = 0.5$.
\end{enumerate}
For each rule, we compute $\alpha_{1/2}(p(\theta))$ as $\theta$
varies.

The results show no distinguishing property of the Born rule.  The
Born-rule $\alpha_{1/2}$ curve has mean $0.910$ and spread $0.069$;
the squared-Born curve has mean $0.916$ and spread $0.069$; the
absolute-value curve has mean $0.898$ and spread $0.041$; the
uniform rule trivially gives a single point at $\alpha_{1/2} =
0.890$.  No rule produces a constant $\alpha_{1/2}$, no rule
produces an extremum that other rules lack, and the Born rule is
not distinguished in any measurable way.

\subsection{Qutrit confirmation}

We repeat the analysis for $d = 3$ (qutrits) with body
probabilities $(p, (1-p)/2, (1-p)/2)$ and model tracking
$\alpha$ on the diagonal.  The $d = 3$ half-saturation curve also
varies with $p$, with mean $\alpha_{1/2} = 0.848$ and spread
$0.023$.  The qualitative behaviour matches $d = 2$: no special
Born-rule property is observed.  The result is not a $d = 2$
artefact.

\medskip
\noindent\textbf{Test A verdict:} The static half-saturation
conjecture is falsified.  The tracking accuracy $\alpha_{1/2}(p)$
varies with $p$, but Born-rule body probabilities show no
distinguished behaviour at half saturation.

% ══════════════════════════════════════════════════════════════════
\section{Test B: Lindblad Dynamics}\label{sec:testB}

\subsection{Pilot trajectory}

Before the full parameter sweep, we examine a single trajectory
with $\theta = \pi/4$ (maximally entangled), $\gamma_D = 1.0$,
$g = 1.0$.  At $t = 0$, the state is pure with
$\IvN = 2\ln 2 = 1.386\nats$ and $\SvN(B) = \ln 2 =
0.693\nats$, giving $r = 2$ and $\rhoQ = -1.386\nats$.

As the system evolves, $\SvN(BM)$ increases (the state becomes
mixed), reducing $\IvN$.  Meanwhile $\SvN(B)$ remains approximately
constant at $\ln 2$ (the Bell state has a maximally mixed reduced
state, and local dephasing does not change the body's marginal
entropy).  The ratio $r$ decreases monotonically from $2$ toward
$1$.

Crucially, $r$ does \emph{not} pass through $r = 1/2$.  It
decreases from $2$ to $1$ and stops, because the exchange
Hamiltonian drives the model to perfectly track the body:
$\IvN \to \SvN(B)$ as $t \to \infty$.  Throughout the entire
trajectory, $r \ge 1$, so $\rhoQ \le 0$ at all times.  The
trajectory functional $\muQ = 0$ identically.

\subsection{Symmetric sweep: 1\,200 trajectories}

We sweep over $100$ values of $\theta \in (0.05, \pi/2 - 0.05)$,
$3$ dephasing rates $\gamma_D \in \{0.1, 1.0, 10.0\}$, and $4$
coupling strengths $g \in \{0.1, 0.5, 1.0, 2.0\}$, yielding
$1\,200$ Lindblad trajectories.  For every trajectory:
\begin{itemize}
  \item $\IvN(t)/\SvN(B)(t) \ge 1.000$ at all stored time steps.
  \item $\rhoQ(t) \le 0$ at all stored time steps.
  \item $\muQ(\theta) = 0$ to within $10^{-15}\nats\!\cdot\!
    \mathrm{time}$.
\end{itemize}
The ratio $r(t)$ decreases monotonically from $2$ (pure entangled
initial state) to $1$ (perfect tracking steady state) for all
parameter combinations.  No trajectory enters the under-correlated
regime $r < 1$ where $\rhoQ > 0$.

The minimum value of $r$ across all $1\,200$ trajectories is
$1.0000$ (to numerical precision $10^{-15}$).  This is not a
marginal result: the ratio approaches $1$ from above and remains
there, with no oscillation or undershoot.

\subsection{Asymmetric extension: 700 trajectories}

To test whether the symmetric dephasing is responsible for the
trivial result, we run $700$ additional trajectories with
asymmetric dephasing.  We fix $g = 1.0$ and $\gamma_0 + \gamma_1
= 2.0$, varying the ratio $\gamma_0/\gamma_1 \in \{0.1, 0.2,
0.5, 1.0, 2.0, 5.0, 10.0\}$ across $100$ values of $\theta$.

The result is identical: $\rhoQ(t) \le 0$ throughout, $\muQ = 0$
identically for all $\gamma_0/\gamma_1$ ratios and all $\theta$.
Asymmetric dephasing does not create an under-correlated transient.

\subsection{Qutrit spot check}

Three spot-check trajectories at $d = 3$ with $\theta \in
\{\pi/8, \pi/4, 3\pi/8\}$, using the generalised exchange
(swap) Hamiltonian and projective dephasing, confirm the same
behaviour: $r \in [1, 2]$, $\rhoQ \le 0$, $\muQ = 0$.  The
maximum $\rhoQ$ across the $d = 3$ trajectories is
$-9.6 \times 10^{-10}\nats$---negative to numerical precision.
The falsification is not an artefact of the qubit ($d = 2$)
Hilbert space.

\medskip
\noindent\textbf{Test B verdict:} The dynamical Born--Fisher--Experiential
conjecture is falsified.  The trajectory functional $\muQ(\theta)
\equiv 0$ for all initial states and all tested decoherence
parameters.  No selection mechanism exists in this model class.

% ══════════════════════════════════════════════════════════════════
\section{Physical Mechanism}\label{sec:mechanism}

The falsification has a clear physical origin.  The exchange
Hamiltonian \eqref{eq:exchange} generates a ``tracking'' interaction
that continuously swaps information from $B$ to $M$.  Combined with
dephasing on $B$, the dynamics proceeds as follows.

At $t = 0$, the pure entangled state $\ket{\psi(\theta)}$ has $r =
2$: the subsystems share twice as much mutual information as the
body's marginal entropy allows classically.  This is the
over-correlated regime, where the excess correlation is due to
entanglement.  In particular, the excess mutual information
$\IvN - \SvN(B)$ for these entangled states includes a quantum
discord contribution~\cite{OllivierZurek2001} that has no
classical analogue.

As dephasing destroys the off-diagonal coherences of $\rho_{BM}$,
the joint entropy $\SvN(BM)$ increases, reducing $\IvN = \SvN(B) +
\SvN(M) - \SvN(BM)$.  Simultaneously, the exchange Hamiltonian
rebuilds classical correlations between $B$ and $M$.

The net effect is that the system transitions \emph{directly} from
the over-correlated regime ($r > 1$, entangled) to the
perfectly-correlated regime ($r = 1$, $\IvN = \SvN(B)$).  It never
passes through an under-correlated phase ($r < 1$) where $\rhoQ$
would be positive.  The trajectory in $(r, t)$ space is a monotone
decreasing curve from $r = 2$ to $r = 1$, parameterised by the
competition between dephasing (which destroys quantum correlations)
and exchange (which builds classical correlations).

This mechanism is specific to the exchange-plus-dephasing model.
The exchange Hamiltonian is ``too good'' at tracking: it ensures
$\IvN \ge \SvN(B)$ at all times, preventing the system from
entering the under-correlated regime.  A weaker coupling mechanism
or a different decoherence channel might produce temporary
under-correlation ($r < 1$), which is the regime where the
experiential density conjecture would have non-trivial content.

Quantitatively, the off-diagonal elements of the joint density
matrix decay as $|\rho_{BM}[0,3](t)| = |\rho_{BM}[0,3](0)|\,
e^{-2\gamma_D t}$ under pure dephasing (validated to relative error
$5 \times 10^{-11}$), while the exchange coupling preserves the
diagonal correlations.  At late times $t \gg \tau_D$, the state is
approximately diagonal with $\IvN \approx \SvN(B)$, giving
$\rhoQ \approx 0^{-}$ (approaching zero from below).

The numerical observation $r(t) \ge 1$ for all trajectories admits
an analytical proof for the model class considered here.

\begin{proposition}[Analytical bound $r \ge 1$ for
exchange-plus-dephasing]\label{prop:r-bound}
  For initial states of the form
  $\ket{\psi(\theta)} = \cos\theta\ket{00} + \sin\theta\ket{11}$
  evolving under the Lindblad equation~\eqref{eq:lindblad} with
  exchange Hamiltonian~\eqref{eq:exchange} and
  dephasing~\eqref{eq:collapse}, the information ratio satisfies
  $r(t) \ge 1$ for all $t \ge 0$.
\end{proposition}

\begin{proof}
The initial state has support only in the $\{\ket{00},\ket{11}\}$
subspace, so the initial density matrix $\rho_{BM}(0)$ has non-zero
entries only at positions $(0,0)$, $(0,3)$, $(3,0)$, and $(3,3)$ in
the computational basis.  We show that the diagonal populations are
time-independent, which fixes the marginal entropies and forces
$r \ge 1$.

\emph{Step 1: The exchange Hamiltonian preserves the populations.}
The exchange Hamiltonian $H = g(\sigma_x \otimes \sigma_x +
\sigma_y \otimes \sigma_y)$ satisfies $H\ket{00} = 0$ and
$H\ket{11} = 0$: explicitly, $(\sigma_x \otimes \sigma_x)\ket{00}
= \ket{11}$ while $(\sigma_y \otimes \sigma_y)\ket{00} =
-\ket{11}$, so $H\ket{00} = g(\ket{11} - \ket{11}) = 0$, and
similarly for $\ket{11}$.  Therefore the commutator $[H, \rho]$
contributes only to off-diagonal coherences and does not change
the diagonal entries of $\rho_{BM}$.

\emph{Step 2: Dephasing preserves the populations.}
The collapse operator $L = \sqrt{\gamma_D}\,(\sigma_z \otimes I)$
is diagonal in the computational basis: $L\ket{ij} = \pm
\sqrt{\gamma_D}\ket{ij}$.  The Lindblad dissipator $L\rho
L^\dagger - \tfrac{1}{2}\{L^\dagger L, \rho\}$ therefore preserves
all diagonal matrix elements of $\rho_{BM}$ while damping
off-diagonal coherences.

\emph{Step 3: Marginal entropies are constant.}
Since the populations $\rho_{BM}[k,k]$ are time-independent, the
partial traces $\rho_B = \Tr_M(\rho_{BM})$ and $\rho_M =
\Tr_B(\rho_{BM})$ are constant.  Hence $\SvN(B)$ and $\SvN(M)$
are constant, with $\SvN(B) = \SvN(M) \equiv S_0$.

\emph{Step 4: Joint entropy increases monotonically.}
The initial pure state has $\SvN(BM) = 0$.  Under the Lindblad
dynamics, dephasing increases $\SvN(BM)$ monotonically from $0$
toward the entropy of the fully diagonal state.  Since the
populations are $(\cos^2\theta, 0, 0, \sin^2\theta)$, the
late-time diagonal state has $\SvN(BM) \to
H_{\mathrm{bin}}(\cos^2\theta) = S_0$.  Therefore
$\SvN(BM)(t) \in [0, S_0]$ for all $t$.

\emph{Step 5: The ratio $r$ is bounded below by $1$.}
Using $\IvN = \SvN(B) + \SvN(M) - \SvN(BM) = 2S_0 - \SvN(BM)$:
\begin{equation}\label{eq:r-bound}
  r(t) = \frac{2S_0 - \SvN(BM)(t)}{S_0}
       = 2 - \frac{\SvN(BM)(t)}{S_0}
       \ge 2 - 1 = 1,
\end{equation}
since $\SvN(BM)(t) \le S_0$.  The ratio decreases monotonically
from $r(0) = 2$ toward $r(\infty) = 1$.
\end{proof}

\begin{remark}
This result extends beyond the specific initial states considered
here.  Any Lindblad dynamics that preserves the diagonal populations
of $\rho_{BM}$ in the computational basis while increasing
$\SvN(BM)$ will maintain $r(t) \ge 1$ provided the initial state
satisfies $r(0) \ge 1$.  The essential mechanism is that the
marginals $\rho_B$ and $\rho_M$ are frozen, so $\IvN$ can only
decrease as the joint state decoheres---but not below $\SvN(B)$.
\end{remark}

% ══════════════════════════════════════════════════════════════════
\section{Discussion}\label{sec:discussion}

\subsection{Status of the conjecture}

The Born--Fisher--Experiential conjecture
(\cref{conj:BFE}) is cleanly falsified for the qubit toy model
class studied here.  The falsification is twofold:
\begin{enumerate}
  \item \textbf{Static (Test A):} In the fully decohered diagonal-state
    regime, the half-saturation tracking accuracy $\alpha_{1/2}(p)$
    varies with $p$ but Born-rule body probabilities are
    indistinguishable from non-Born alternatives.
  \item \textbf{Dynamic (Test B):} Under Lindblad evolution with
    exchange coupling and dephasing, the experiential density is
    non-positive throughout the evolution, making the trajectory
    functional $\muQ$ identically zero with no selection mechanism.
\end{enumerate}

\subsection{What is not falsified}

The experiential density functional $\rhoQ$ itself is well-defined
and well-behaved.  It correctly identifies the parabolic peak
structure ($\rhoQ^{\max} = \SvN(B)/4$ at $r = 1/2$) in the static
diagonal-state regime.  The issue is not with the functional but
with the conjecture that Born-rule probabilities are selected by
it.  In the dynamical setting, the system never enters the regime
where $\rhoQ > 0$, so the functional has no opportunity to act as a
selector.

The pure-state negativity $\rhoQ < 0$ for all entangled pure states
(\cref{eq:pure-ratio}) is a structural property of the naive quantum
extension, not specific to the Born--Fisher conjecture.  Any future
quantum application of the experiential density would need to either
modify the functional form---for example, replacing $\IvN$ with a
quantity bounded by $\SvN(B)$ such as the Holevo quantity or
accessible information---or restrict the domain to decohered
(mixed) states where $r \le 1$.  This structural limitation
motivates the search for alternative quantum extensions of the
experiential density.

Structurally, the experiential density is parallel to Tononi's
integrated information $\Phi$~\cite{Tononi2004}: both are bipartite
information functionals that quantify meaningful internal structure.
Quantum extensions of integrated information
theory~\cite{Zanardi2018} face analogous challenges with entangled
states, where quantum correlations inflate the information measures
beyond their classically meaningful ranges.

The Lipschitz stability of the classical experiential
measure~\cite{Ehrlich2026b} is also unaffected.  The classical
functional $\rho = I(B;M)(1 - I/H(B))$ retains its properties;
the quantum extension simply reveals that the half-saturation
condition is not reached during Markovian Lindblad evolution with
exchange coupling.

\subsection{Future directions}

The key question is whether there exist quantum channels for which
$r = \IvN/\SvN(B)$ transiently dips below $1$, creating a positive
$\rhoQ$ window.  Several avenues merit investigation:

\begin{itemize}
  \item \textbf{Non-Markovian channels.}  Memory effects can
    produce oscillations in $\IvN(t)$ that might temporarily reduce
    $r$ below $1$.  Nakajima--Zwanzig or process-tensor approaches
    could test this.
  \item \textbf{Amplitude damping.}  Unlike dephasing, amplitude
    damping changes the populations of $\rho_B$, potentially
    creating transient under-correlation.
  \item \textbf{Weak coupling.}  The exchange Hamiltonian produces
    perfect tracking ($r \to 1$) at late times.  A weaker or
    imperfect coupling that allows $r < 1$ could open the
    experiential-density window.
  \item \textbf{Higher dimensions.}  While the $d = 3$ spot check
    shows the same qualitative behaviour, larger Hilbert spaces may
    exhibit richer dynamics.
\end{itemize}

The connection to the Baez--Dolan groupoid
cardinality~\cite{Baez2001}, which motivates the
$1/|\mathrm{Aut}(x)|$ weighting on the structure space, remains
conceptually interesting even though the specific half-saturation
conjecture fails.  The experiential density framework may find its
quantum application in a different variational principle.

% ══════════════════════════════════════════════════════════════════
\section{Conclusion}\label{sec:conclusion}

We have tested the Born--Fisher--Experiential conjecture using a
qubit toy model with two complementary approaches: static
diagonal-state analysis and full Lindblad master equation dynamics.
Both tests yield clear negative results.

The static test shows that the half-saturation tracking accuracy
$\alpha_{1/2}(p)$ varies with body probability $p$ (spread $0.069$
for $d = 2$), but Born-rule probabilities display no special
behaviour compared with non-Born alternatives.

The dynamical test, spanning $1\,900^{+}$ Lindblad trajectories
with symmetric dephasing, asymmetric dephasing, and qutrit
extensions, shows that the experiential density $\rhoQ(t) \le 0$
throughout the entire evolution.  The information ratio
$\IvN/\SvN(B)$ monotonically decreases from $2$ to $1$, never
entering the under-correlated regime where $\rhoQ > 0$.  The
trajectory functional $\muQ(\theta) = 0$ identically (to
$10^{-15}\nats\!\cdot\!\mathrm{time}$), with no dependence on
initial state, dephasing rate, or coupling strength.

The physical mechanism is the exchange Hamiltonian's
perfect-tracking behaviour: the system transitions directly from
over-correlated (entangled) to perfectly correlated (classical
tracking) without passing through an under-correlated phase.  The
conjecture is cleanly falsified for this model class.  Non-Markovian
channels and amplitude damping remain open avenues for modified
conjectures.

% ══════════════════════════════════════════════════════════════════
\appendix
\section{Numerical Validation}\label{app:validation}

All computations were implemented in Python (NumPy) without
external quantum libraries.  We report two suites of validation
checks: 10 quantum-information sanity checks for the core utilities
and 6 Lindblad integrator validation checks.

\subsection{Quantum information sanity checks}

\begin{table}[H]
\centering
\small
\begin{tabular}{clcc}
  \toprule
  \# & Check & Expected & Result \\
  \midrule
  (a) & $\Tr(\rho) = 1$, PSD for $\rho_{BM}(0.5, 0.75)$
      & pass & pass \\
  (b) & Pure Bell state $\IvN = 2\ln 2$
      & $1.38629\nats$ & $1.38629\nats$ ($\delta < 10^{-12}$) \\
  (c) & Product state $\IvN = 0$
      & $0$ & $< 10^{-12}$ \\
  (d) & $\SvN(I/2) = \ln 2$
      & $0.69315\nats$ & $0.69315\nats$ ($\delta < 10^{-12}$) \\
  (e) & $\SvN(I/3) = \ln 3$
      & $1.09861\nats$ & $1.09861\nats$ ($\delta < 10^{-12}$) \\
  (f) & Araki--Lieb bound, 10 random states
      & $0 \le \IvN \le 2\min(\SvN(B), \SvN(M))$ & 10/10 pass \\
  (g) & Symmetry $I(p,\alpha) = I(1-p,\alpha)$
      & $\delta < 10^{-12}$ & max err $< 2 \times 10^{-16}$ \\
  (h) & Eigenvalue cutoff stability ($\varepsilon = 10^{-15,-14,-13}$)
      & agreement & $\delta < 10^{-13}$ \\
  (i) & Perfect tracking $\alpha = 1$: $r = 1$
      & $1.0$ & $1.0$ ($\delta < 10^{-12}$) \\
  (j) & No tracking $\alpha = 0.5$: $\IvN = 0$
      & $0$ & $< 10^{-12}$ \\
  \bottomrule
\end{tabular}
\caption{Quantum information utility validation.
All 10 checks pass to machine precision or better.}
\label{tab:qi-checks}
\end{table}

\subsection{Lindblad integrator validation}

\begin{table}[H]
\centering
\small
\begin{tabular}{clcc}
  \toprule
  \# & Check & Tolerance & Result \\
  \midrule
  (a) & Trace preservation: $\max|\Tr(\rho) - 1|$
      & $10^{-10}$ & $< 10^{-16}$ \\
  (b) & Positivity: $\min(\lambda_i)$
      & $\ge -10^{-10}$ & $\ge -2 \times 10^{-31}$ \\
  (c) & Pure dephasing rate: $|\rho_{BM}[0,3](t)|$
      & $5\%$ relative & $5 \times 10^{-11}$ relative \\
  (d) & Initial condition: $\IvN(0) = 2\ln 2$
      & $10^{-12}$ & $\delta < 10^{-12}$ \\
  (e) & Late-time diagonal: off-diagonal sum
      & $< 10^{-4}$ & $< 2 \times 10^{-9}$ \\
  (f) & RK4 convergence: $|\muQ(n{=}2000) - \muQ(n{=}4000)|$
      & $10^{-4}$ & $< 10^{-15}$ \\
  \bottomrule
\end{tabular}
\caption{Lindblad integrator validation.
All 6 checks pass with margins of several orders of magnitude.
The trace is preserved to $10^{-16}$, positivity to $10^{-31}$,
and RK4 convergence is verified to $O(h^4)$.}
\label{tab:lindblad-checks}
\end{table}

% ══════════════════════════════════════════════════════════════════

\begin{thebibliography}{99}

\bibitem{Baez2001}
J.~C.~Baez and J.~Dolan,
``From finite sets to Feynman diagrams,''
in \textit{Mathematics Unlimited---2001 and Beyond},
B.~Engquist and W.~Schmid, eds.
(Springer, Berlin, 2001), pp.~29--50.
\texttt{arXiv:math/0004133}.

\bibitem{Valentini2005}
A.~Valentini and H.~Westman,
``Dynamical origin of quantum probabilities,''
\textit{Proc.\ R.\ Soc.\ A} \textbf{461}, 253--272 (2005).

\bibitem{Gleason1957}
A.~M.~Gleason,
``Measures on the closed subspaces of a Hilbert space,''
\textit{J.\ Math.\ Mech.} \textbf{6}, 885--893 (1957).

\bibitem{Chiribella2011}
G.~Chiribella, G.~M.~D'Ariano, and P.~Perinotti,
``Informational derivation of quantum theory,''
\textit{Phys.\ Rev.\ A} \textbf{84}, 012311 (2011).

\bibitem{NielsenChuang}
M.~A.~Nielsen and I.~L.~Chuang,
\textit{Quantum Computation and Quantum Information}
(Cambridge University Press, Cambridge, 2010), 10th anniversary ed.

\bibitem{BreuerPetruccione}
H.-P.~Breuer and F.~Petruccione,
\textit{The Theory of Open Quantum Systems}
(Oxford University Press, Oxford, 2002).

\bibitem{Towler2012}
M.~D.~Towler, N.~J.~Russell, and A.~Valentini,
``Time scales for dynamical relaxation to the Born rule,''
\textit{Proc.\ R.\ Soc.\ A} \textbf{468}, 990--1013 (2012).

\bibitem{Ehrlich2026a}
B.~Ehrlich,
``Exponential suppression of transient-basin contributions in
trajectory-weighted Markov chain measures,''
preprint (2026).

\bibitem{Ehrlich2026b}
B.~Ehrlich,
``Lipschitz Stability of the Experiential Density Functional,''
preprint (2026).

\bibitem{Deutsch1999}
D.~Deutsch,
``Quantum theory of probability and decisions,''
\textit{Proc.\ R.\ Soc.\ A} \textbf{455}, 3129--3137 (1999).

\bibitem{Wallace2007}
D.~Wallace,
``Quantum probability from subjunctive reasoning,''
\textit{Stud.\ Hist.\ Phil.\ Mod.\ Phys.} \textbf{38}, 311--332 (2007).

\bibitem{Zurek2005}
W.~H.~Zurek,
``Probabilities from entanglement, Born's rule $p_k = |\psi_k|^2$
from envariance,''
\textit{Phys.\ Rev.\ A} \textbf{71}, 052105 (2005).

\bibitem{OllivierZurek2001}
H.~Ollivier and W.~H.~Zurek,
``Quantum discord: a measure of the quantumness of correlations,''
\textit{Phys.\ Rev.\ Lett.} \textbf{88}, 017901 (2001).

\bibitem{Tononi2004}
G.~Tononi,
``An information integration theory of consciousness,''
\textit{BMC Neuroscience} \textbf{5}, 42 (2004).

\bibitem{Zanardi2018}
P.~Zanardi, G.~Styliaris, and L.~Venuti,
``Measures of coherence-generating power for quantum unital
operations,''
\textit{Phys.\ Rev.\ A} \textbf{95}, 052307 (2017).

\end{thebibliography}

\end{document}
